\documentclass[12pt]{article} 
\usepackage{longtable}
\usepackage{makecell}
\usepackage{float}
\usepackage{graphicx}
\usepackage{bm}
\usepackage{placeins}
\usepackage{booktabs}
\usepackage{gensymb}
\usepackage{amssymb}
\usepackage{indentfirst}
\usepackage{threeparttable} 
\usepackage{multirow}
\usepackage{subfigure} % 导言区引入
\usepackage{tabularx}
\usepackage{aligned-overset}
\usepackage[slantfont,boldfont]{xeCJK}
\usepackage{fontspec}
\renewcommand{\arraystretch}{1.5}
% \usepackage[paperwidth=21cm,paperheight=29.7cm]{geometry}
\usepackage[a4paper,margin=2cm]{geometry}
\setCJKmainfont{SimSun}
\setmainfont{SimSun}
\setsansfont{SimSun}
% 设置正文字体为四号（12pt）
% \renewcommand{\CJKmainfont}{SimSun}
\renewcommand{\CJKfamilydefault}{\CJKrmdefault}
\renewcommand{\normalsize}{\fontsize{12pt}{18pt}\selectfont}
\normalsize

\title{迈克尔逊干涉仪实验报告}
\author{2411545 邱凯锐}
\date{2025.11.14}

\begin{document}
\maketitle

\section{实验目的}
1. 了解迈克尔逊干涉仪的结构原理并掌握调节方法 

2. 观察等厚干涉、等倾干涉

3.  测量钠双线的平均波长和波长差

\section{实验仪器}
迈克尔逊干涉仪  钠灯

\section{实验原理}

\subsection{迈克尔逊干涉仪}

迈克尔逊干涉仪是一个分振幅法的双光束干涉仪，其光路如 Figure 1 所
示，它由反射镜 $M_1$、$M_2$、分束镜 $P_1$ 和补偿板 $P_2$ 组成。其中 $M_1$ 是一个
固定反射镜，反射镜 $M_2$ 可以沿光轴前后移动，它们分别放置在两个相
互垂直臂中；分束镜和
补偿板与两个反射镜均
成 $45^\circ$ ，且相互平行；分
束镜 $P_1$ 的一个面镀有半
透半反膜，它能将入射
光等强度地分为两束；
补偿板是一个与分束镜
厚度和折射率完全相同
的玻璃板。

光源发出的光经分
束镜被分成等强度的两
束光 1 和 2 ，光束 1 和 2 分别经反射镜 $M_1$ 和 $M_2$ 反射后，再次经分光镜
$P_1$ 向 $E$ 处传播。由于光束 2 在传播过程中三次穿过分束镜，而光束 1
只有一次穿过分束镜。由于玻璃存在色散，不同波长的光在干涉仪中具
不同的光程差，为此，在反射镜 $M_1$ 和反射镜之间加入一个补偿板，这
样光线 1 同样在相同的玻璃板中穿过三次，使所有波长的光可以同时获
得接近零的光程差，这对于实现白光的干涉是绝对必要的前提。在单色
光入射时，补偿板可以两臂的光程达到完全对称，如果没用补偿板，单
色光的干涉条纹将变的很复杂。 

在迈克尔逊干涉仪上可以实现等倾和等厚两种干涉。
对于等倾干涉，如果反射镜 $M_1$ 和 $M_2$ 到分束镜的距离差为 $d$ ，光线入射角为 $i$ 时，两反射镜反射光的光程差 $\Delta$ 为：
\begin{equation}
    \Delta=2nd\cos i'=2d\sqrt{n^2-\sin^2i}
\end{equation}
其中，$n$ 为两臂中介质的折射率，$i$ 和 $i'$ 分别为光线入射到 $M_2$ 和 $M_1$ 上
的入射角，当迈克尔逊干涉仪的两臂中介质相同时，$i=i'$ 。

\begin{figure}[H]
    \centering
    \includegraphics[width=0.65\textwidth]{fig1.png}
    \caption{迈克尔逊干涉仪光路}
\end{figure}

\begin{figure}[H]
    \centering
    \includegraphics[width=0.65\textwidth]{fig2.png}
    \caption{等倾干涉的光程差}
\end{figure}

当两臂中介质的折射率一定，且 $d$ 不变时，光程差 $\Delta$ 只取决于入射角
$i$ ，在 $E$ 处观察时，对于相同入射角的光，形成一个以光轴为中心的园环。
当：
\begin{equation}
    \begin{aligned}
        &\Delta=k\lambda,\quad &k=0,1,2\cdots\quad &\text{为亮条纹}\\
        &\Delta=\frac{2k+1}{2}\lambda \quad &k=0,1,2,\cdots\quad &\text{为暗条纹}
    \end{aligned}
\end{equation} 
由此可以看出，迈克尔逊干涉仪中，等倾干涉条纹级次是中间大边上小。
当改变 $d$，光程差也相应发生改变，这时在干涉条纹中心会出现“冒出”
和“缩进”的现象。当 $d$ 增加 $\lambda/2$，相应的光程差增加 $\lambda$ ，在中心的条纹
干涉级次由 $k$ 变为 $k+1$，这样就会“冒出”一个条纹；当 $d$ 减少 $\lambda/2$ ，相
应的光程差减少 $\lambda$ ，在中心的条纹干涉级次由 $k$ 变为 $k-1$ ，这样就会“缩
进”一个条纹。因此，根据“冒出”或“缩进”条纹的个数可以确定 $d$
的改变量，它可以用来进行长度测量，其精度是波长量级。当“冒出”
或“缩进”了 $N$ 个条纹，$d$ 的改变量 $\Delta d$ 为： 
\begin{equation}
    \Delta d=N\frac \lambda 2
\end{equation}

\begin{figure}[H]
    \centering
    \includegraphics[width=0.65\textwidth]{fig3.png}
    \caption{等厚干涉的光程差}
\end{figure}

当 $M'_1$、$M_2$ 不相互平行，即 $M_1$ 和 $M_2$
之间有一小的夹角 $\theta$，光束 1 和 2，在反射后
相遇在 $M'_2$ 附近的 $P$ 点，如 Figure 3 所示，即
等厚干涉条纹定为于 $M'_2$ 的表面，两束光
的光程差仍然是：$\Delta=2d\cos i$
，如果采用
孔径角不太大的观察系统或者眼睛观察，
$i'$ 一般是很小的，这样可以将光程差式(1)
展开到 $i'$ 的二次项
\begin{equation}
    \Delta=2d\left(1-\frac {i'^2}{2} \right) 
\end{equation}
当 $d$ 很小，并满足以下时
\begin{equation}
    di'^2<<\lambda
\end{equation}
 $i'$ 对光程差的影响可以忽略不计，式(4)可以写成 
\begin{equation}
    \Delta=2d
\end{equation}
 这时具有相同 $d$ 值的反射光束有相同的光程差，形成平行于交棱的干涉
直条纹。但如果 $d$ 较大，不满足式(5)时，随着 $i'$ 的增大干涉条纹会发生
弯曲。

\subsection{利用迈克尔逊干涉仪测量钠黄光双线的平均波长}
在等倾干涉的条件下，由式（3）可知，在实验中，通过调节 $M_1$ 和 $M_2'$ 的间距 $d$，记录干涉条纹“冒出”（或“缩进”）的条纹数 $N$，可以得到钠黄光双线的平均波长
\begin{equation}
    \bar \lambda=\frac{2\Delta d}{N}
\end{equation} 

\subsection{利用迈克尔逊干涉仪测量钠黄光双线的波长差}
钠黄光是有两个强度基本相同、波长相差不大的双线组成，它们之
间的波长差可以利用迈克尔逊干涉仪测得。由式(1)可知，对于不同波长
的光，即使 $d$ 相同，同一干涉级次的干涉条纹也具有不同的角距离 $i'$ ，
它们不相重合。

对于钠黄光双线来说，它们的波长分别为 $\lambda_1$、$\lambda_2$，彼此甚为相近。
如果定义 $\lambda_2>\lambda_1$
，可能出现这样的情况，当 $d = d_1$ 时，$\lambda_2$ 的 $k_2$ 级亮条纹
与 $\lambda_1$ 的 $k_1$ 级的暗条纹重合，导致视场的条纹对比度为零，这时有 
\begin{equation}
    2d_1\cos i'=k_2\lambda_2=(k_1+\frac 12)\lambda_!
\end{equation}
即，$\lambda_2$ 干涉级次改变了 $\delta k$、$\lambda_1$ 的干涉级次改变了 $\delta k+1/2$。这时视场中
两个波长的亮条纹和亮条纹重合，暗条纹和暗条纹重合，干涉条纹非常
清晰，其对比度为 1。继续改变 $d$， $d = d_3$，使视场中的干涉条纹可见都
再次变为零，这时由 $d_1$ 到 $d_3$，$\lambda_2$ 干涉级次改变了 $\delta k'$、$\lambda_1$ 的干涉级次改
变了 $\delta k'＋1$，于是有：
\begin{equation}
    2d_3\cos i'=(k_2+\delta k)\lambda_2=(k_1+\delta k'+\frac 32)\lambda_1
\end{equation} 
对于中心的条纹有 $i'\approx 0$,由式（7）和（9）相减可得
\begin{equation}
    \delta \lambda=\lambda_1/\delta k'
\end{equation}
且，
\begin{equation}
    \delta d=d_3-d_1=\delta k'\lambda_2=(\delta k'+1)\lambda_1
\end{equation}
将式（11）代入式（10），并考虑到 $\bar\lambda =\frac 12(\lambda_1+\lambda_2)>>\delta \lambda$ 可得
\begin{equation}
    \delta \lambda=\frac{\lambda_1\lambda_2}{\delta d}=\frac{\bar\lambda^2}{\delta d}
\end{equation}
因此，通过记录两次条纹可见都为零之间，干涉级次的变为 $\delta k'$、或者反射镜 $M_2$ 移动的距离 $\delta d$,分别通过式（10）和（12）得到 $\delta \lambda$。

\section{仪器介绍}
实验中所使用的WSM－100迈克尔逊干涉仪，其结构如 Figure 4 所示。

固定反射镜10，它倾角可以通过在它背面的三个调整螺丝进行粗调，
细调时，分别用1和4通过弹簧改变在水平和垂直方向上的倾角。 

可移动反射镜 9，它的倾角调整也是通过其背面的三个螺丝进行，
但它没有细调。可移动反射镜是安装在一个滑块上，通过丝杠 16 的转
动可以使滑块沿导轨前后移动。手轮2是通过一个传动比为2:1齿轮组
带动丝杠转动，反射镜的移动距离可以在直尺 17 上读到毫米，在读数
窗11 中读到0.01，微动手轮4是通过一个100:1涡轮蜗杆减速机构来实
现反射镜的精细移动，在微动手轮上的最小读数值为0.0001mm。

分束镜18和补偿板12，以及固定反射镜10是安装在一个固定架上，
这个固定架又与导轨固定在一起，使得整个光路处在一个比较稳定的状
态。

导轨是固定在底座7上，底座由三个底角6支撑，可以通过调整螺
丝调节整个仪器的水平。

\begin{figure}[H]
    \centering
    \includegraphics[width=0.6\textwidth]{fig4.png}
    \caption{迈克尔逊干涉仪}
\end{figure}

\section{实验内容}
\subsection{调节迈克尔逊干涉仪}

1.首先调整可移动反射镜 $M_2$ 到分束镜之间的距离，使之与固定反射
镜 $M_1$ 到分束镜之间的距离大致相等。原因是 $d$ 是过大，不利于观
察到干涉条纹。

2.调整反射镜 $M_1$ 和 $M_2$ 倾斜，使 $M_1⊥M_2$，即使 $M'_1∥M_2$。具体方法
是：打开钠灯，在钠灯和分束镜之间加入一个画有一个黑色十字的
毛玻璃，这时的光源为面光源，在视场中可以看到两组共 4 个十字
的象，其中两个较深，另两个较浅。调节 $M_1$ 后面的三个螺丝，使
两组象完全重合，这时两个反射镜接近垂直，应该可以看到干涉条
纹。

如果没有干涉条纹出现，可重复上面两个步骤。

\subsection{观察等倾和等厚干涉现象}
观察等倾干涉：观察在看到条纹后，利用手轮 2 和微动手轮 3 改变
光程差，使干涉条纹疏密适中；并调节 $M_1$ 后面的三个螺丝使圆环处在
视场中心。这时上、下、左、右移动眼睛，改变观察方向。当眼睛上下
移动时，干涉条纹有“冒出”和“缩进”，调节反射镜 $M_1$ 下方的弹簧螺
丝 4，精细改变 $M_1$ 在垂直方向上的倾斜；当眼睛左右移动时，干涉条纹
有“冒出”和“缩进”，调节反射镜 $M_1$ 侧面的弹簧螺丝 1，精细改变
$M_1$ 在水平方向上的倾斜；直到在改变观察方向时，干涉条纹没有“冒出”
和“缩进”，且圆心随眼睛的移动而移动，干涉条纹的大小不发生改变，
此时，观察到的干涉条纹时严格的等倾干涉条纹。

观察等厚干涉：在观察到等倾干涉条纹后，改变反射镜 $M_2$ 的位置，
逐步减小光程差，在 $d≈0$ 的附近，利用弹簧螺丝 1 或弹簧螺丝 4 改变
反射镜 $M_1$ 的倾斜，使 $M'_1$ 和 $M_2$ 形成一个微小的夹角，这时在视场中会
出现略有弯曲的等厚干涉条纹。

\begin{figure}[H]
    \centering
    \begin{minipage}{0.45\textwidth}
        \centering
        \includegraphics[width=\textwidth]{fig2.jpg}
        \caption{(a) 等倾干涉条纹}
    \end{minipage}
    \hfill
    \begin{minipage}{0.45\textwidth}
        \centering
        \includegraphics[width=\textwidth]{fig1.jpg}
        \caption{(b) 等厚干涉条纹}
    \end{minipage}
\end{figure}

\subsection{测量钠黄光的平均波长和钠双线的波长差}
记录微调手轮 3 每转一周，$\Delta d =10\mu m$，等倾干涉条纹“冒出”或“缩
进”个数，根据式（3）计算钠黄光的平均波长，并与钠黄光的实际平均波
长 589.3nm 进行比较，给出相对误差，分析误差的来源。
\begin{table}[H]
    \centering
    \caption{测量钠黄光的平均波长}
    \vspace{2pt} 
    \begin{minipage}{0.7\textwidth}
            \raggedleft
            $\Delta d=10\ \mathrm{\mu m}$
    \end{minipage}\\[2pt]
    \begin{tabular}{cccc}
        \hline\hline
        次数 $i$ & 1 & 2 & 3  \\ 
        $N$ & 33 & 31 & 34 \\ 
        \hline\hline
    \end{tabular}
\end{table}
计算得到
\begin{equation}
    \bar\lambda=\frac{2\Delta d}{\bar N}=612.24\ \mathrm{nm}
\end{equation}
相对误差
\begin{equation}
    E=\left|\frac{\lambda-\lambda_0}{\lambda_0}\right|\times 100\%=3.9\%
\end{equation}
误差来源:

1.条纹数 $N$ 只能取整数值，且难以准确数清；

2.反射镜 $M_1$ 和 $M_2$ 不可能调节至完全平行且难以使得 $i'\approx 0$ ；

3.分束镜和反射镜表明并非完美平整，存在弯曲导致位移偏差。

记录相邻两次视场中干涉条纹可见度为零时，反射镜所移动的距离，
利用式（12）计算钠双线的波长差。
\begin{table}[H]
    \centering
    \caption{测量钠双线的波长差}
    \vspace{2pt} 
    \begin{minipage}{0.7\textwidth}
            \raggedleft
            $\bar \lambda=612.24\ \mathrm{nm}$
    \end{minipage}\\[2pt]
    \begin{tabular}{cccc}
        \hline\hline
        次数 $i$ & 1 & 2 & 3  \\ 
        $\delta d\ \mathrm{mm}$ & 0.27772 & 0.30556 & 0.29300 \\ 
        \hline\hline
    \end{tabular}
\end{table}
计算得到
\begin{equation}
    \delta \lambda=1.28\ \mathrm{nm}
\end{equation}

\section{思考题}

1.\textbf{根据迈克耳孙干涉仪的光路，说明各光学元件的作用。}

（1）分光镜：把入射光分成两束强度近似相等的光；把两臂反射回来的光重新合束，使它们叠加形成干涉条纹。

（2）反射镜：各自反射由分光镜分出的光束，使其沿原路返回；通过调整 $M_1$ 或 $M_2$ 的位置（微调螺钉），改变两臂光程差，从而使干涉条纹

（3）补偿板：补偿两束光在玻璃中走过的厚度差，使两臂光程匹配；补偿色散，使得两束光可以出现清晰明亮的干涉条纹。

2.\textbf{在迈克耳孙干涉仪上看到的等倾圆条纹与牛顿环实验中的干涉圆
条纹有哪些区别？实验中如何避免回空差？}

区别：牛顿环本质上是等厚干涉，且中心暗斑对应着 0 级暗条纹。

如何避免回空差：（1）调节螺旋时保持一个方向；（2）微调镜时不要快速跳动，缓慢调整，使条纹运动均匀，不要大幅来回转动。

3.\textbf{在观察非定域干涉时，为什么当 $d$ 足够大时，屏上看不到干涉条
纹了?}

光源非严格单色光，持续发光时间 $\tau_0$ 有限，存在时间上的相干性，对应着相干长度 $L_0$。如果 $d$ 过大，就看不到干涉条纹了。

4.\textbf{在实验中有时会观察到椭圆或者马鞍型的条纹，思考成因。} 

等倾条纹椭圆形：（1）两反射镜 $M_1$、$M_2$ 不完全垂直于光路；（2）分光镜或补偿板存在楔角且不完全平行。

等厚条纹马鞍型：（1）反射镜表面不平；（2）组件不在同一平面，光路出现像差；

5.\textbf{改用白色台灯作光源会有什么现象？}

（1）出现彩色的干涉条纹；（2）衬比度下降；（3）条纹范围非常小；

6.\textbf{如果用激光作光源，迈克耳孙干涉仪的两臂不相等，对现象有影
响吗？}

因为激光的相干性较高，所以（1）两臂不等长时，仍然能看到干涉条纹；（2）两臂差别太大时，条纹亮度会减弱或消失

\end{document}